\documentclass[12pt,a4paper]{article}

\usepackage[utf8]{inputenc}
\usepackage[T1]{fontenc}
\usepackage[english]{babel}
\usepackage{lmodern}
\usepackage{microtype}
\usepackage[a4paper,margin=2.5cm]{geometry}
\usepackage{setspace}
\onehalfspacing

\usepackage{amsmath,amssymb,amsfonts,mathtools}
\usepackage{amsthm}
\usepackage{booktabs}
\usepackage{xcolor}
\usepackage[hidelinks]{hyperref}
\usepackage[round,authoryear]{natbib}

\newtheorem{remark}{Remark}

\newcommand{\E}{\mathbb{E}}
\newcommand{\tD}{\widetilde{D}}
\newcommand{\tZ}{\widetilde{Z}}
\newcommand{\tY}{\widetilde{Y}}
\newcommand{\hattheta}{\hat{\theta}}
\newcommand{\SE}{\mathrm{SE}}

\title{\textbf{Why IV Estimators Have Larger Standard Errors:\\
Mechanics, Implications, and Evidence from the Application}}
\author{}
\date{}

\begin{document}
\maketitle
\thispagestyle{empty}

\section*{Overview}

A consistent pattern across all specifications in this thesis is that IV-based estimators
--- PLR-IV and Wald-AIPW --- produce substantially larger standard errors and wider
confidence intervals than their non-IV counterparts PLR and AIPW\@.  This note explains
the mechanical reason for this precision loss, derives the variance inflation formally,
and interprets the specific pattern observed in the Card (1995) application.

%----------------------------------------------------------------------
\section{The Fundamental Mechanic: Division by the First Stage}
%----------------------------------------------------------------------

Every IV estimator can be written as a ratio of two quantities: a reduced-form
relationship between the instrument and the outcome, and a first-stage relationship
between the instrument and the treatment.  In the simplest case of the plain Wald
estimator this is
\begin{align}
  \hat{\tau}_{\mathrm{Wald}}
  \;=\;
  \frac{
    \E_N[Y_i \mid Z_i=1] - \E_N[Y_i \mid Z_i=0]
  }{
    \E_N[D_i \mid Z_i=1] - \E_N[D_i \mid Z_i=0]
  }
  \;=\;
  \frac{\hat\pi_{\mathrm{RF}}}{\hat\pi_{\mathrm{FS}}}.
  \label{eq:wald_ratio}
\end{align}
The variance of a ratio estimator follows by the delta method.  Let
$\sigma^2_{\mathrm{RF}} = \mathrm{Var}(\hat\pi_{\mathrm{RF}})$ and
$\sigma^2_{\mathrm{FS}} = \mathrm{Var}(\hat\pi_{\mathrm{FS}})$.  Then
\begin{align}
  \mathrm{Var}(\hat\tau_{\mathrm{Wald}})
  \;\approx\;
  \frac{\sigma^2_{\mathrm{RF}}}{\pi_{\mathrm{FS}}^2}
  \;+\;
  \frac{\pi_{\mathrm{RF}}^2 \,\sigma^2_{\mathrm{FS}}}{\pi_{\mathrm{FS}}^4}.
  \label{eq:delta_method}
\end{align}
The key term is $1/\pi_{\mathrm{FS}}^2$.  Whenever the instrument explains only a small
share of the variation in $D_i$, the first stage $\pi_{\mathrm{FS}}$ is small, this
amplification factor is large, and all sampling uncertainty in the numerator is
magnified into the final estimate.  This is the \emph{variance inflation} inherent to
IV\@.

In the DML setting, the same logic applies after residualization.  The PLR-IV estimator
solves the moment condition
$\E[(\tY_i - \theta_0 \tD_i)\tZ_i] = 0$
and has the closed-form solution
\begin{align}
  \hattheta_{\mathrm{PLR\text{-}IV}}
  \;=\;
  \frac{\sum_i \hat{V}_i \,\tY_i}{\sum_i \hat{V}_i \,\tD_i},
  \qquad
  \hat{V}_i = Z_i - \hat{m}(X_i).
  \label{eq:plriv_ratio}
\end{align}
The denominator $\sum_i \hat{V}_i \tD_i$ is the residualized first stage --- the
covariance between the instrument residual and the treatment residual after partialling
out all covariate information.  If the instrument is only weakly related to treatment
after controlling for $X_i$, this denominator is small and variance inflates by exactly
the same mechanism as in \eqref{eq:delta_method}.

%----------------------------------------------------------------------
\section{Why OLS and AIPW-ATE Are More Precise}
%----------------------------------------------------------------------

OLS and AIPW-ATE do not divide by a first stage.  They use \emph{all} variation in
$D_i$ to identify the treatment effect.  In the residualized PLR representation, for
instance, the estimator is
\begin{align}
  \hattheta_{\mathrm{PLR}}
  \;=\;
  \frac{\sum_i \tD_i \,\tY_i}{\sum_i \tD_i^2},
\end{align}
and the denominator $\sum_i \tD_i^2$ is the full residual variance of the treatment ---
always large, irrespective of whether an instrument is available.

The fundamental trade-off is therefore between \emph{consistency} and
\emph{precision}.  OLS/AIPW achieves high precision by using all treatment variation,
but at the cost of consistency: when $D_i$ is endogenous, these estimators are biased.
IV achieves consistency by discarding all endogenous variation and relying only on
instrument-driven variation, but that variation is typically a small fraction of the
total, so precision falls.  This is not a failure of IV; it is the price of causal
identification under endogeneity.

\begin{remark}[Subpopulation interpretation of precision loss]
An equivalent way to understand the precision loss is through the effective sample size
(ESS).  IV estimators implicitly work with a \emph{sub-experiment}: the variation in
outcomes induced by the instrument among compliers.  If the instrument affects treatment
for, say, 15--20\% of the sample, the IV estimator is effectively drawing on that
fraction of the data for identification.  Formally, the ESS of the IV outcome weights
$\omega_i$ is $(\sum_i \omega_i^2)^{-1}$, which collapses when weights are concentrated
on a small subgroup.  Low ESS $\Rightarrow$ high variance.
\end{remark}

%----------------------------------------------------------------------
\section{Three Compounding Factors in the Card Application}
%----------------------------------------------------------------------

In the Card (1995) application, three features compound the general variance inflation
and produce the very large standard errors observed across all IV specifications.

\medskip\noindent\textbf{(i) Moderate instrument strength.}
The instrument is a binary indicator for whether an individual grew up near a four-year
college (\texttt{nearc4}).  While the first-stage relationship in Card (1995) is
statistically significant, it is not strong by the standards of modern applied
econometrics.  A first-stage coefficient of roughly 0.15--0.20 means the denominator
in \eqref{eq:plriv_ratio} is small, and the amplification factor $1/\hat\pi_{\mathrm{FS}}^2$
is correspondingly large.

\medskip\noindent\textbf{(ii) Binary instrument and binary treatment.}
Both $Z_i$ (college proximity) and the treatment variables ($D_1 = \mathbf{1}\{\text{educ}>12\}$,
$D_2 = \mathbf{1}\{\text{educ}\geq 16\}$) are binary.  Binary variables have limited
variance relative to continuous ones, which mechanically limits the first-stage signal
and reduces the denominator in the IV ratio further.

\medskip\noindent\textbf{(iii) Compliers are a small subpopulation.}
Not everyone responds to the college proximity instrument.  The complier share ---
those who would attend more education if and only if they grew up near a college ---
is a fraction of the sample.  The LATE estimator is identified solely off this group.
This is reflected in the effective sample sizes reported in the weight diagnostics:
$\mathrm{ESS} \approx 2$--$3$ for PLR-IV and Wald-AIPW across all specifications,
despite $N = 3{,}010$ observations.  Estimating a causal effect from the information
content of two to three effectively independent units produces wide confidence intervals
almost by construction.

%----------------------------------------------------------------------
\section{Empirical Pattern in the Results}
%----------------------------------------------------------------------

Table~\ref{tab:se_comparison} collects the estimates and standard errors from the four
main DML specifications to illustrate the precision gap.

\begin{table}[htbp]
\centering
\caption{Estimates and standard errors by estimator and specification}
\label{tab:se_comparison}
\smallskip
\begin{tabular}{llcccc}
\toprule
 & & \multicolumn{2}{c}{Card specification}
   & \multicolumn{2}{c}{Kitagawa specification} \\
\cmidrule(lr){3-4}\cmidrule(lr){5-6}
Estimator & Target & $D_1$ (somecol) & $D_2$ (educ16)
          & $D_1$ (somecol) & $D_2$ (educ16) \\
\midrule
PLR       & $\theta_0$ (structural)
  & 0.239 (0.018) & 0.320 (0.019) & 0.113 (0.015) & 0.152 (0.018) \\
AIPW-ATE  & ATE
  & ---           & 0.342 (0.015) & 0.117 (0.015) & 0.172 (0.018) \\
\midrule
PLR-IV    & $\theta_0$ (structural, IV)
  & 0.681 (0.289) & 1.140 (0.535) & 0.593 (0.340) & 1.030 (0.684) \\
Wald-AIPW & LATE
  & 0.279 (0.214) & 0.530 (0.398) & 0.292 (0.249) & 0.535 (0.462) \\
\bottomrule
\end{tabular}
\smallskip\\
\raggedright\small
\textit{Notes:} Standard errors in parentheses.  $D_1 = \mathbf{1}\{\mathrm{educ}>12\}$,
$D_2 = \mathbf{1}\{\mathrm{educ}\geq 16\}$.  The Card specification includes experience,
experience squared, regional indicators, and race/location dummies.  The Kitagawa
specification includes race and location dummies only.  DML estimates use 5-fold
cross-fitting with tuned random forests.
\end{table}

\noindent Several patterns are immediately apparent.

\medskip\noindent\textbf{SE ratios of 10--25$\times$.}
Comparing IV to non-IV standard errors within the same specification, the ratio ranges
from roughly 12$\times$ (PLR vs.\ PLR-IV, Card/$D_1$: $0.018$ vs.\ $0.289$) to 26$\times$
(AIPW-ATE vs.\ Wald-AIPW, Kitagawa/$D_2$: $0.018$ vs.\ $0.462$).  This is not
a marginal difference in precision; it is an order-of-magnitude change.

\medskip\noindent\textbf{Significance disappears for IV estimators.}
All non-IV estimators (PLR, AIPW-ATE, AIPW-ATT, AIPW-ATU) are significant at the
0.1\% level across all specifications.  Among IV estimators, Wald-AIPW is not
statistically significant in any specification.  PLR-IV reaches significance only in
the Card specification (both $D_1$ and $D_2$) but not in the Kitagawa specification.
This divergence in significance is entirely explained by the standard error gap; the
point estimates are not zero and are economically meaningful.

\medskip\noindent\textbf{IV point estimates are larger, not smaller.}
A natural concern when IV standard errors are large is that the estimates themselves are
implausible.  Here the opposite is true: IV estimates are consistently \emph{larger}
than OLS/AIPW counterparts.  This is consistent with a downward ability bias in OLS ---
individuals with higher ability sort into more education and also earn higher wages, but
if ability is unobserved, OLS conflates the education effect with the ability selection,
potentially biasing estimates \emph{upward} in the standard Mincer setup.  However, Card
(1995) and subsequent work have argued that the return to education for compliers
(near-college individuals induced into more schooling) may genuinely exceed the average,
because compliers are on a margin where the education investment is particularly valuable.
Both explanations are consistent with larger IV than OLS estimates; the data alone cannot
distinguish them.

\medskip\noindent\textbf{Wider confidence intervals but directionally consistent.}
Despite lacking statistical significance, the Wald-AIPW confidence intervals are wide
but centred on positive values that are directionally consistent across all four
specifications (Card/$D_1$: 0.28; Card/$D_2$: 0.53; Kitagawa/$D_1$: 0.29;
Kitagawa/$D_2$: 0.54).  The 95\% confidence interval for Wald-AIPW under
Kitagawa/$D_2$, for instance, runs approximately $[-0.37,\, 1.44]$ --- it is wide, but
it does not provide evidence \emph{against} a positive return to college completion.

%----------------------------------------------------------------------
\section{Implication for the Empirical Comparison}
%----------------------------------------------------------------------

The large standard errors of IV estimators should not be interpreted as a methodological
failure.  They honestly reflect the limited first-stage variation available in the
Card (1995) design.  The comparison between IV and non-IV estimators remains
informative for two reasons.

First, the point estimates suggest an upward correction when moving from OLS to IV,
which is consistent with the theoretical prediction that IV corrects for endogeneity
bias in the returns to education.  That the correction cannot be precisely estimated
with this instrument is a limitation of the data and instrument, not of the estimators.

Second, the precision comparison itself is a finding.  It illustrates that obtaining
clean causal estimates in observational settings with endogenous treatment requires
strong instruments.  The college proximity instrument of Card (1995) is historically
important and widely used, but it provides limited statistical power for precisely
estimating the LATE in the modern DML framework with flexible nuisance estimation.
This motivates reporting both the IV and non-IV estimates jointly: together they bound
the plausible range of causal effects while making the endogeneity concern and its cost
in precision transparent.

%----------------------------------------------------------------------
\bibliographystyle{apalike}
\bibliography{references}

\end{document}
